% Generated draft: requires mathematical and source review.
\documentclass[11pt]{article}
\usepackage[margin=25mm]{geometry}
\usepackage{fontspec}
\setmainfont{Latin Modern Roman}
\usepackage{amsmath,amssymb,mathtools,unicode-math}
\setmathfont{Latin Modern Math}
\usepackage{enumitem,needspace,xurl,hyperref}
\setlength{\parindent}{0pt}
\setlength{\parskip}{5pt}
\setlength{\emergencystretch}{3em}
\newcommand{\N}{\mathbb{N}}
\newcommand{\Z}{\mathbb{Z}}
\newcommand{\Q}{\mathbb{Q}}
\newcommand{\R}{\mathbb{R}}
\newcommand{\C}{\mathbb{C}}
\newcommand{\F}{\mathbb{F}}
\newcommand{\A}{\mathbb{A}}
\newcommand{\PP}{\mathbb{P}}
\newcommand{\E}{\mathbb{E}}
\newcommand{\eps}{\varepsilon}
\newcommand{\dd}{\,\mathrm{d}}
\begin{document}
\section*{UnsolvedMath ID: 30004283}
\textbf{Source problem code:} catalogue entry
\subsection{Definitions and mathematical statement}
Let \(R\) be the hyperfinite factor of type \(\mathrm{II}_1\). Interpret the category in the source as the rigid \(C^*\)-tensor category \(\mathcal C_R\) of dualizable Hilbert \(R\)-\(R\) bimodules: a bimodule is a Hilbert space with commuting normal unital actions of \(R\) and \(R^{\mathrm{op}}\), and morphisms are bounded bimodule intertwiners. The tensor product is Connes fusion \(H\otimes_R K\), and the tensor unit is \(L^2(R)\). Dualizability means admitting a conjugate bimodule and evaluation and coevaluation morphisms satisfying the duality identities; these are the bifinite bimodules referred to in the supplied literature assessment. This specifies the standard operator-algebraic interpretation, which should be checked against the cited report.

Let \(\operatorname{Irr}(\mathcal C_R)\) denote the isomorphism classes of nonzero simple objects, where simplicity means \(\operatorname{End}_{\mathcal C_R}(X)=\C\). For a group \(G\), a grading is an assignment \(d:\operatorname{Irr}(\mathcal C_R)\to G\) such that \(d(L^2(R))=e_G\) and
\[
d(Z)=d(X)d(Y)
\quad\text{whenever}\quad
\operatorname{Hom}_{\mathcal C_R}(Z,X\otimes_R Y)\ne 0
\]
for simple \(X,Y,Z\). Thus every simple constituent of \(X\otimes_R Y\) has degree \(d(X)d(Y)\).

The universal grading group \(U_R\) is equipped with such a grading \(u\) for which every grading \(d\) into any group \(G\) factors uniquely as \(d=\phi\circ u\) through a group homomorphism \(\phi:U_R\to G\). Equivalently, its defining presentation is
\[
U_R=\left\langle
u_X\ (X\in\operatorname{Irr}(\mathcal C_R))
\ \middle|\ 
u_{L^2(R)}=1,\quad
u_Z=u_Xu_Y\text{ if }\operatorname{Hom}_{\mathcal C_R}(Z,X\otimes_R Y)\ne0
\right\rangle,
\]
where each symbol \(u_X\) depends only on the isomorphism class of \(X\).

\textbf{Problem.} Determine \(U_R\) up to group isomorphism, together with its universal grading of \(\mathcal C_R\).

The supplied dated assessment states that positive real modular distortion gives a grading and that its universality has not been proved. This is supplied literature curation, not an independently verified assertion here. The record gives neither a normalization of modular distortion nor a precise conjectural identification of \(U_R\); no such identification is imposed on the problem.
\subsection{Short English statement}
Determine the group that universally records degrees of simple dualizable bimodules over the hyperfinite \(\mathrm{II}_1\) factor, with multiplication compatible with every simple constituent of Connes fusion. The supplied record labels the problem open; its dated literature assessment and the precise categorical conventions require independent review.
\subsection{Sources}
\begin{itemize}
\item Supplied local UnsolvedMath registry record, permanent ID \(30004283\), title ``Universal Grading of Hyperfinite Factor Bimodules,'' source code \texttt{catalogue entry}, in the Oberwolfach Reports Open Problems collection. The supplied proposed year is 2019; no individual proposer is named. The record labels the problem \texttt{open} and its statement \texttt{corrected\_verified}, explaining the latter as editorial repair with mathematical content preserved. These are supplied registry labels, not independent verification.
\item Supplied report citation: \emph{Subfactors and Applications} (2020), Oberwolfach Reports, DOI 10.4171/owr/2019/49, pp.~3053--3140. \url{https://doi.org/10.4171/owr/2019/49}. The background describes these reports as documenting mathematical workshops and open-problem discussions.
\item Supplied literature source: Penneys, ``unitary pivotal structures and modular distortion.'' \url{https://dml.cz/bitstream/handle/10338.dmlcz/153425/HigherStructures_004-2020-2_2.pdf}.
\item Supplied literature source: Bischoff et al., ``distortion for multifactor bimodules.'' \url{https://doi.org/10.4171/DM/1011}.
\item Supplied literature source: ``original Oberwolfach report.'' \url{https://doi.org/10.4171/owr/2019/49}.
\item Literature assessment checked at \texttt{2026-08-22}: the supplied curation labels the problem open, says that positive real modular distortion gives a grading, and says that its universality is not proved. This date and assessment are supplied dated curation derived from the OWR record, not independently verified literature findings. Time-sensitive status and cited claims require independent verification.
\end{itemize}
\end{document}
